%%%%%%%%%%%%%%%%%%%%%%%%%%%%%%%%%%%%%%%%%%%%%%%%%%%%%%%%%%%%%%%%%%%%%%%%%%%%%%%%
%%  EINLAGE 6: QUANTITATIVE AUSSAGE AUS EINEM ANDEREN VORTRAG
%%  vorgeschlagene Abgabe: 21.02.2027 (freiwillig)
%%  Struktur: Beschreibung (1/4–1/2 S.) + Diskussion der quantitativen Physik (1–3 S.) + Reflexion (1–2 S.)
%%  Vortrag/Folien der anderen Gruppe in literatur.bib eintragen (Vorlage VORLAGE-Vortrag).
%%%%%%%%%%%%%%%%%%%%%%%%%%%%%%%%%%%%%%%%%%%%%%%%%%%%%%%%%%%%%%%%%%%%%%%%%%%%%%%%
% \FremdVortragKey wird in einstellungen.tex gesetzt

\section{Einlage: Eine quantitative wissenschaftliche Aussage aus einem anderen Vortrag, die für Ihr Projekt eine Relevanz hat oder anwendbar ist.}
\label{sec:fremdvortrag}

\begin{hinweise}
Die Quelle dieser Einlage oder die Art der Einlage ist aus den von Ihnen gehörten Vorträgen (ausser dem
eigenen) frei wählbar. Sie muss mit Ihrem Schwerpunktthema einen Zusammenhang haben.

\smallskip\emph{Kandidaten aus der Themenliste WS 26/27 mit Bezug zu \Schwerpunkt:}
\enquote{Stabilität von rollenden Fahrrädern} (Stabilität aus linearisierten Bewegungsgleichungen,
Eigenwerte), \enquote{Untersuchung eines Doppelpendels} (numerische Integration, Chaos -- vgl.\ chaotisches
Verhalten bei schneller Drehung in \cite{\PaperKey}), \enquote{Dynamik eines rotierenden Tischtennisballs}
(Kreiseldynamik). Erst nach den Vorträgen am 28.\,01.\,2027 festlegen.
\end{hinweise}

\subsection{Beschreibung}
\begin{hinweise}
Die Beschreibung gibt alle notwendigen Informationen an, die die Einlage beschreiben, also zum Beispiel die
Gruppe, die den Vortrag gehalten hat, das Thema dieser Gruppe und eine kurze Inhaltsübersicht, aber keine
Nacherzählung (min. $1/4$ Seite, max. $1/2$ Seite!).
\end{hinweise}
\begin{umfang}{Einlage 6 (fremder Vortrag): Beschreibung}{0.25}{0.5}
\platzhalter{Gruppe, Thema, Datum, gewählte Aussage (Folie \dots) \cite{\FremdVortragKey}; kurze
Inhaltsübersicht; Folie(n) in Anhang~\ref{anh:fremdvortrag}}
\end{umfang}

\subsection{Diskussion der quantitativen Physik (oder der Mathematik)}
\begin{hinweise}
Wenn Sie für diese Einlage eine mathematische oder physikalische Aussage gewählt haben, beschreiben Sie hier
die Physik oder Mathematik dieser Einlage.
\begin{itemize}\itemsep0pt
\item Was ist die Hauptaussage?
\item Was sind die Voraussetzungen?
\item Wie hilft physikalisch die Physik (Mathematik oder Wissenschaft) dieser Einlage bei Ihrem Schwerpunktthema?
\item Welche quantitativen Aussagen unterlegt mit physikalischen Gleichungen sind für Sie besonders wichtig?
\item Wo gab es Schwierigkeiten beim Verständnis der Physik (Mathematik oder Wissenschaft)? Welche?
\item Wie können Sie mit dieser Aussage Ihr Projekt und Ihre Behandlung des Schwerpunktthemas verbessern oder korrigieren? Welchen Einfluss hätte dies?
\end{itemize}
\textbf{Die Diskussion der Physik sollte mindestens eine Seite und maximal 3 Seiten lang sein. Die
Längenangaben sind ohne Titel, Bilder, Tabellen und Tabellen- oder Bildlegenden und ohne übergrosse Abstände
zu Gleichungen zu verstehen.}
\end{hinweise}
\begin{umfang}{Einlage 6 (fremder Vortrag): Diskussion der Physik}{1}{3}
\platzhalter{Hauptaussage, Voraussetzungen, Gleichungen (nummeriert), Anwendung auf \Schwerpunkt,
quantitative Abschätzung des Einflusses auf das eigene Projekt}
\end{umfang}

\subsection{Reflexion}
\begin{hinweise}
In der Reflexion sollen Sie sich Gedanken machen über die Wirkung dieser Einlage und der Arbeit mit ihr auf
Ihren Lernfortschritt.
\begin{itemize}\itemsep0pt
\item In welchem Zusammenhang steht die Einlage mit Ihrem Projektthema?
\item In welchem Zusammenhang steht die Einlage mit Ihrem Schwerpunktthema?
\item Was haben Sie bei der Beschäftigung mit dieser Einlage gelernt?
\item Wie hat Ihnen die Einlage beim Verständnis Ihres Schwerpunktthemas geholfen?
\item Wo gab es Schwierigkeiten bei der Beschäftigung mit der Einlage?
\item Was hätten Sie anders gemacht, wenn Sie die Einlage erstellt hätten?
\item Was sagen Ihre Kommilitoninnen oder Kommilitonen dazu, als Kommentar oder wenn Sie mit Ihnen über die Einlage diskutieren?
\item Enthält die die verwendete Aussage logische oder physikalische Fehler?
\item Gab es bei der Beschäftigung mit der Einlage eine überraschende neue Erkenntnis?
\end{itemize}
\textbf{Die Reflexion für eine Einlage sollte mindestens eine Seite und maximal 2 Seiten lang sein. Die
Längenangaben sind ohne Titel, Bilder, Tabellen und Tabellen- oder Bildlegenden zu verstehen.}
\end{hinweise}
\begin{umfang}{Einlage 6 (fremder Vortrag): Reflexion}{1}{2}
\platzhalter{Reflexion in Fließtext -- Leitfragen erkennbar beantworten}
\end{umfang}

%\subsection{Eigene Bewertungskriterien}
%...
